
% ==================================================================================================================================
% Analyse Synatxique (Grammaires LL(k))

\section{Analyse Syntaxique et approche descendante (Grammaires $LL(k)$)}

Lors de la dérivation d'un mot par une grammaire, on peut se heurter à un problème de taille : quel règle de dérivation choisir 
pour obtenir un arbre de dérivation acceptant ? 
Nous allons donc introduire un type de grammaire nous permettant de régler ce problème. Un seul type de dérivation 
sera possible et nous saurons exactement combien de caractères regarder dans notre mot pour choisir la règle 
de dérivation à appliquer. 

Les grammaires $LL(k)$ vont nous permettre de faire directement le lien avec le processus de \emph{parsing}, indispensable 
en compilation. 

\subsection{Définition}

\begin{definition}[Grammaire $LL(k)$]
    Soit $k \in \N$ et $G$ une grammaire. On dit que $G$ est $LL(k)$ si on peut analyser n’importe quel mot 
    de son langage de gauche à droite, en construisant la dérivation la plus à gauche, et en regardant au plus $k$ 
    symboles d’entrée à la fois, pour décider de manière univoque quelle production appliquer à chaque étape.
\end{definition}

\begin{remark}
    Détaillons un peu cette définition. La grammaire $G$ est $LL(k)$ si elle satsifait trois critères 
    principaux : 
    \begin{itemize}
        \item \textbf{L (Left to right)} : Lorsque l'on dérive un mot par la grammaire, on veut pouvoir 
        le lire de gauche à droite en commençant par le premier symbole. 
        \item \textbf{L (Leftmost derivation)} : Lorsque l'on remplace une variable non terminale 
        de l'espression en cours de dérivation, on veut que ce soit celle la plus à gauche. 
        \item \textbf{k (k lookahead)} : $k$ est le nombre de symbole que le parser peut regarder 
        pour déterminer quelle règle utiliser sans les consommer. 
    \end{itemize}
\end{remark}

Intuitivement, on cherche une grammaire prédictive. Pour chaque variable à chaque étape de la dérivation, 
le parser doit savoir exactement quelle règle appliquer en regardant $k$ lettres. Évidement, on ne considèrera 
pas de grammaire ambiguë. Nous allons donc développer un algorithme pour construire une \emph{table de dérivation} 
qui nous donnera la règle à appliquer à chaque étape de l'analyse syntaxique. 

\subsection{Ensemble First, Follow, $k$-lookahead}

\begin{definition}[$k$ - Lookadhead]
    Soit $G = ( \Sigma, V, S, P)$ une grammaire $LL(k)$ et une règle de production $A \rightarrow \alpha$. On définit 
    l'ensemble $LA(k, A \rightarrow \alpha)$ comme l'ensemble des chaînes de longueur $ \leqslant k$ qui déclenchent 
    la règle de production $A \rightarrow \alpha$ tel que : 
        \[ \boxed{LA(k, A \rightarrow \alpha) := \{w \in \Sigma^{ \leqslant k} | w \text{ peut apparaître si on choisit la règle } A \rightarrow \alpha\}} \] 
\end{definition}

$LA(k, A \rightarrow \alpha)$ est donc l'ensemble des lettres que doit regarder le parser pour déterminer 
quelle règle de production appliquer. 
Son calcul permet aussi de vérifier qu'un grammaire est $LL(k)$. En effet si les $k$-lookahead des règles 
d'un même non terminal ont des mots en communs, la grammaire n'est pas $LL(k)$ (i.e on ne saura pas choisir 
quelle règle appliquer). 

\begin{definition}[Ensemble First]
    Soit $G = ( \Sigma, V, S, P)$ une grammaire $LL(k)$. 
    Soit $X \in V$ un état \emph{non terminal}. On définit l'ensemble $First(X)$ comme l'ensemble 
    des \emph{élément terminaux de $ \Sigma$} qui peuvent apparaître en premier lors d'une dérivation de $X$. 
    De plus, si $X$ peut produire un $ \varepsilon$, alors $ \varepsilon \in  First(X)$. Plus formellement : 
        \[ First_k( \alpha) := 
            \left\{
                u \in \Sigma^*, \quad 
                \begin{cases}
                    \text{ si } |u| < k \text{ et } \alpha \to^* u \\ 
                    \text{ si } |u| = k \text{ et } \alpha \to^* u \beta \text{ avec } \beta \in (V \cup \Sigma)^*      
                \end{cases}
            \right\}
        \] 
\end{definition}

Cet ensemble permet de déterminer l'ensemble des mots de longueur inférieure ou égale à $k$ qui peuvent 
dériver du mot auquel on l'applique. En général on commence le calcul à partir d'un non-terminal.

\begin{prop}[Règles de la calcul (First)]
    Soit $G = ( \Sigma, V, S, P)$ une grammaire $LL(k)$, on a les règles de calcul suivantes : 
    \begin{itemize}
        \item Soit $a \in \Sigma$ un terminal, on a : $ First(a) = \{a\}$ 
        \item $First( \varepsilon) = \{ \varepsilon\}$ 
        \item Soit $A \in V$ une variable avec les productions : $A \rightarrow \alpha_1 | \dots | \alpha_n$ on a :
            \[ First(A) = \bigcup_{i \in \llbracket 1, n \rrbracket} First( \alpha_i) \] 
    \end{itemize}
\end{prop}

\begin{example}[Ensemble First]
    Soit $G$ la grammaire aux règles de production suivantes : 
        $$ 
            \begin{cases}
                S' \to S \$ \\ 
                S \to aAe | bc | \varepsilon \\ 
                A \to d | \varepsilon 
            \end{cases}
        $$
    On a :
        $$ First_3(S) = First_3(aAe) \cup First_3(bc) \cap First_3( \varepsilon) $$ 
    Or : 
        $$ First_3(aAe) = a \cdot First_2(A) \cdot e = a \{d, \varepsilon\} e = \{ade, ae\} $$ 
    D'où : 
        $$ First_3(S) = \{ade, ae, bd, \varepsilon\} $$
\end{example}

\begin{definition}[Ensemble Follow]
    Soit $G = ( \Sigma, V, S, P)$ une grammaire $LL(k)$. 
    Soit $X \in V$ un état \emph{non terminal} et $k \in \N$. On définit $Follow_k(A)$ comme l'ensemble 
    des mots de longueur exactement $k$ qui peuvent suivre le non terminal $A$. 
\end{definition}

\begin{prop}[Règle de calcul (Follow)]
    Soit $G = ( \Sigma, V, S, P)$ une grammaire $LL(k)$, on a alors : 
        $$ Follow_k(A) = \bigcup_{Y \to \gamma A \delta} First_k( \delta Follow_k(Y)) $$
    Si $\varepsilon \in First_k(\delta)$, on ajoute aussi $Follow_k(Y)$.
\end{prop}

\begin{example}[Ensemble Follow]
    Soit $G$ la grammaire aux règles de production suivantes : 
        $$ 
            \begin{cases}
                S' \to S \$ \\ 
                S \to aA | bABc \\ 
                A \to d \\ 
                B \to e | \varepsilon 
            \end{cases}
        $$
    On remarque que : 
        \begin{itemize}
            \item $First_3(Bc) = First_3(B) \cdot c = \{ec, c\}$ 
            \item $Follow_3(S) = First_3(Follow_3(\$^3)) = \{\$^3\}$
        \end{itemize}
        \begin{align*}
            Follow_3(A) &= First_3(Follow_3(S)) \cup First_3(Bc Follow_3(S)) \\ 
            &= \{\$^3\} \cup First_3(Bc \cdot Follow_3(S)) \\ 
            &= \{\$^3\} \cup \{ec \cdot First_1(Follow_1(S))\} \cup \{c \cdot First_2(Follow_2(S))\} \\ 
            &= \{\$^3\} \cup \{ec\$\} \cup \{c\$^2\} \\ 
            &= \{\$^3, ec\$, c\$^2\} 
        \end{align*}
\end{example}

\begin{example}[Follow avec récursivité]
    Soit $G$ la grammaire aux règles de production suivantes : 
        $$ 
            \begin{cases}
                S' \to S \$^3 \\ 
                S \to aA \\ 
                A \to bAc | \varepsilon
            \end{cases}
        $$ 
    On a ici : 
    \begin{align*}
        Follow_3(A) &= First_3(Follow_3(S)) \cup First_3(Follow_3(A)) \\ 
        &= First_3(\$^3 \cdot Follow_3(S')) \cup c \cdot First_2(Follow_2(A)) \\ 
        &= \{\$^3\} \cup \{\$^2\} \cup c \cdot First_1(Follow_1(A)) \\ 
        &= \{\$^3, c\$^2, cc\$, c c c\}
    \end{align*}
\end{example}

\subsection{Règles de calcul et table d'analyse}

\begin{theorem}[Calcul de $LA(k)$]
    Soit $G = ( \Sigma, V, S, P)$ une grammaire $LL(k)$ et une règle de production $A \rightarrow \alpha$.
    On a : 
        \[ \boxed{LA(k, A \rightarrow \alpha) = First_k( \alpha) \cup \{x \in Follow_k(A) | \alpha \text{ peut dériver } \varepsilon\}} \] 
    plus précisément : 
    \begin{itemize}
        \item Si $ \alpha$ ne peut pas dériver $ \varepsilon$ : $LA(k, A \rightarrow \alpha) = First_a( \alpha)$ 
        \item Si $ \alpha$ peut dériver $ \varepsilon$ : $ LA(k, A \rightarrow \alpha) = First_a( \alpha) \cup Follow_k(A) $
    \end{itemize}
\end{theorem}

\begin{proposition}
    On peut réduire le théorème suivant à cette simple formule : 
        $$ LA(k, A \to \alpha) = First_k( \alpha Follow_k(A)) $$ 
\end{proposition}

L'ensemble de ces règles de calcul vont nous permettre de définir quelle règle de dérivation appliquer à chaque 
étape de l'analyse syntaxique d'un mot. Nous regrouperons toutes ces règles dans une table appelée la 
\emph{table d'analyse}. 

\begin{definition}[Table d'analyse]
     Soit $G = ( \Sigma, V, S, P)$ une grammaire $LL(k)$. Pour chaque règle de 
     production $A \longrightarrow \alpha$ et chaque $w \in LA(k, A \longrightarrow \alpha)$, on définit : 
        \[ M[A, w] = \alpha \] 
    Autrement dit, $ M[A,w]$ contient la partie droite de la règle de production à utiliser quand le non-terminal 
    courant est $A$ et que les $k$ prochains symboles de l'entrée forment $w$. 
\end{definition}

La table d’analyse d’une grammaire $LL(k)$ permet, pour chaque non-terminal $A$ et chaque séquence de 
$k$ symboles de lookahead $w$, de déterminer quelle production appliquer.
Pour construire la table d'analyse d'une grammaire $LL(k)$, nous devrons calculer les $k$-lookahead 
pour chaque règle de la grammaire. 

\begin{example}
    Soit $G$ la grammaire suivante : 
        $$ 
            \begin{cases}
                S' \to S\$^3 \\ 
                S \to aAb \\ 
                A \to c | \varepsilon 
            \end{cases}
        $$
    Calculons sa table d'analyse. Pour cela, commençons par le calcul des $3-lookahead$ : 
        \begin{align*}
            LA(3, S' \to S\$^3 &= First_3(S\$^2 \cdots Follow_3(S)) \\ 
            &= First_3(S\$^3) \\
            &= First_3(acb\$^3) \cup First_3(ab\$^3) \\ 
            &= \{acb, ab\$\}
        \end{align*}
        \begin{align*}
            LA(3, S \to aAb) &= First_3(aAb \cdot Follow_3(S)) \\ 
            &= First_3(aAb \cdots First_3(\$^3)) \\ 
            &= First_3(aAb\$^3) \\ 
            &= First_3(acb\$^3) \cup First_3(ab\$^3) \\ 
            &= \{acb, ab\$\}
        \end{align*}
        \begin{align*}
            LA(3, A \to c) &= First_3(c \cdot Follow_3(A)) \\ 
            &= First_3(c \cdot First_3(b \cdot Follow_3(S))) \\ 
            &= First_3(c \cdot First_3(c\$^3)) \\ 
            &= First_3(c \cdot \$^2) = \{cb\$\}
        \end{align*}
        \begin{align*}
            LA(3, A \to \varepsilon) &= First_3(Follow_3(A)) \\ 
            &= First_3(b\$^2) = \{b\$^2\}
        \end{align*}
    
    On peut donc construire la table d'analyse de la grammaire : 

        \[
            \begin{array}{c|cccc}
                \hline 
                \text{Non-terminal} & acb & ab\$ & cb\$ & b\$^2 \\ 
                \hline 
                \hline
                S' & S' \to S\$^3 & S' \to S\$^3 & & \\ 
                S  & S \to aAb  & S \to aAb  & & \\ 
                A  &            &           & A \to c & A \to \varepsilon\\
                \hline 
            \end{array}
        \]
    Chaque cellule indique la production à appliquer si le flux d'entrée commence par la 
    chaîne correspondante de longueur 3.
\end{example}

% \begin{example}
%     Considérons la grammaire :
%         \[
%             \begin{aligned}
%                 S &\to A\,b \mid c \\
%                 A &\to d \mid \varepsilon
%             \end{aligned}
%         \]

%     {Étape 1 : Ensembles FIRST et FOLLOW}
%         \[
%             \begin{array}{lcl}
%                 \text{FIRST}(A) &=& \{ d, \varepsilon \} \\[3pt]
%                 \text{FIRST}(A b) &=& \{ d, b \} \\[3pt]
%                 \text{FIRST}(S) &=& \{ d, b, c \} \\[6pt]
%                 \text{FOLLOW}(S) &=& \{ \$ \} \\[3pt]
%                 \text{FOLLOW}(A) &=& \{ b \}
%             \end{array}
%         \]

%     {Étape 2 : Ensembles LA(1)}

%         \[
%             \begin{array}{|c|c|c|c|}
%                 \hline
%                 \textbf{Production} & \textbf{FIRST} & \textbf{FOLLOW (si $\varepsilon \in \text{FIRST}$)} & \textbf{LA(1)} \\
%                 \hline
%                 S \to A b & \{d, b\} & - & \{d, b\} \\
%                 \hline
%                 S \to c   & \{c\} & - & \{c\} \\
%                 \hline
%                 A \to d   & \{d\} & - & \{d\} \\
%                 \hline
%                 A \to \varepsilon & \{\varepsilon\} & \{b\} & \{b\} \\
%                 \hline
%             \end{array}
%         \]

%    On peut alors écrire la table d'analyse de la grammaire : 

%         \[ 
%             \begin{array}{|c|c|c|c|c|}
%                 \hline 
%                 \textbf{Non Terminal} & \textbf{b} & \textbf{c} & \textbf{d} & \textbf{\$} \\ 
%                 \hline 
%                 {S} & S \rightarrow Ab & S \rightarrow c & S \rightarrow Ab & \text{(erreur)} \\ 
%                 \hline 
%                 {A} & A \rightarrow \varepsilon & \text{(erreur)} & A \rightarrow d & \text{(erreur)} \\ 
%                 \hline
%             \end{array}
%         \] 
% \end{example}

% \begin{example}
%     Considérons la grammaire :
        
%         \[
%             \begin{aligned}
%                 S &\to aA \mid aB \\
%                 A &\to b \\
%                 B &\to c
%             \end{aligned}
%         \]

%     Pour les deux productions de $S$, le premier symbole est $a$ :
%         \[
%             \text{FIRST}(aA) = \text{FIRST}(aB) = \{a\}
%         \]

%     Cette grammaire n'est donc {pas LL(1)} : avec un seul symbole de regard, on ne peut pas distinguer $aA$ de $aB$.

%     {Calculons les ensembles LA(2)} : On regarde deux symboles $(k=2)$ pour décider :

%         \[
%             \begin{array}{|c|c|c|}
%                 \hline
%                 \textbf{Production} & \textbf{Chaînes possibles} & \textbf{LA(2)} \\
%                 \hline
%                 S \to aA & ab & \{ab\} \\
%                 \hline
%                 S \to aB & ac & \{ac\} \\
%                 \hline
%                 A \to b & b & \{b\} \\
%                 \hline
%                 B \to c & c & \{c\} \\
%                 \hline
%             \end{array}
%         \]

%     {Conclusion :} Avec $k=2$, on distingue les deux productions de $S$ :
%     \begin{itemize}
%         \item[-] Si les deux prochains symboles sont \texttt{ab}, on choisit $S \to aA$.
%         \item[-] Si les deux prochains symboles sont \texttt{ac}, on choisit $S \to aB$.
%     \end{itemize}

%         \[
%             \begin{array}{|c|c|c|}
%                 \hline
%                 \textbf{Non-terminal} & \textbf{ab} & \textbf{ac} \\
%                 \hline
%                 S & S \to aA & S \to aB \\
%                 \hline
%             \end{array}
%         \]

%         Pour les non-terminaux $A$ et $B$, un seul symbole de lookahead suffit :
%         \[
%             \begin{array}{|c|c|}
%                 \hline
%                 \textbf{Non-terminal} & \textbf{Symbole} \\
%                 \hline
%                 A & b \Rightarrow A \to b \\
%                 \hline
%                 B & c \Rightarrow B \to c \\
%                 \hline
%             \end{array}
%         \]
%     La grammaire est donc {LL(2)}, mais pas LL(1).
% \end{example}

Ces nouveaux modèles de grammaires, quoique puissants ne permettent pas d'effectuer l'analyse syntaxique 
de n'importe quel langage. En effet, toutes les grammaires ne sont pas $LL(k)$ nous devrions donc 
modifier celles qui ne le sont pas pour effectuer l'analyse. 

% ==================================================================================================================================
% Analyse Syntaxique et approche ascendante (Grammaires LR(k))

\section{Analyse Syntaxique et approche ascendante (Grammaires $LR(k)$)}

\subsection{Principe}

Une autre approche de l'analyse syntaxique grâce aux grammaires et d'utiliser l'approche \emph{descendante}.
En effet, depuis le début nous cherchons à dériver un mot à partir de l'axiome de départ en descendant 
vers des états terminaux. Nous allons ici chercher à faire le contraire : à partir d'un mot nous chercherons 
à le réduire vers l'axiome. C'est ce que l'on appelle l'approche \emph{ascendante}, aussi appelée 
\emph{analyse par décalage/réduction (shift-reduce)}. 

\begin{example}
    Commençons par étudier un exemple pour bien comprendre le principe de ces grammaires. 
    Soit la grammaire $G_1$ suivante : 
        \[ 
            \begin{cases}
                E \to E + E \\ 
                E \to E * E \\ 
                E \to (E) \\ 
                E \to id 
            \end{cases}
        \] 
    Soit $ w \in \Sigma^*$ tel que $w := id + id * id$. Nous cherchons à réduire $w$ grâce aux règles de 
    productions de $G_1$ jusqu'à l'axiome $S$. Résumons tout ça dans un tableau : 
        \[ 
            \begin{array}{|c|c|c|}
                \hline 
                \text{Proto-phrase de droite} & \text{Poignée (handle)} & \text{Production de réduction} \\ 
                \hline 
                id + id * id & id & E \to id \\
                \hline  
                E + id * id & id & E \to id \\ 
                \hline
                E + E * id & id & E \to id \\ 
                \hline 
                E + E * E & E * E & E \to E * E \\ 
                \hline
                E + E & E + E & E \to E + E \\ 
                \hline 
                E & & \\
                \hline 
            \end{array}
        \] 
\end{example}

\subsection{Outils Fondamentaux}

Définissons maintenant quelques outils pour effectuer rigoureusement cette analyse ascendante. 

\begin{definition}[proto-mot]
    Soit $G = (V, \Sigma, S, P)$ une grammaire. Un \emph{proto-mot} (ou \emph{chaîne sentencielle}) est 
    toute chaîne de symboles, terminaux et/ou non terminaux, que l'on peut obtenir en partant de 
    l'axiome et en effectuant un nombre quelconque de dérivations. 
    Formellement, pour la grammaire $G$, on a : 
        \[ \alpha \in (V \cup \Sigma)^* \text{ est un proto-mot si } S \to^* \alpha \] 
\end{definition}

Un proto-mot est simplement une "étape intermédiaire" de la dérivation d'un mot. Dans 
l'ensemble précédent, cela correspond à $ E + E * id$. Il est composé du terminal $id$ et 
du non terminal $E$. 

\begin{definition}[Poignée (Handle)]
    Soient $G = (V, \Sigma, S, P)$ une grammaire et $w$ un proto-mot de $G$. 
    Une \emph{poignée} de $w$ est occurence précise du corps d'une production $A \to \alpha$ dans 
    $w$ telle que : 
        \[ S \to^* \gamma A \delta \to \gamma \alpha \delta = w \]
\end{definition}

Une poignée correspond donc à une sous-chaîne à réduire dans notre mot $w$ à un instant donné de l'analyse ascendante. 
Dans notre exemple précédent, pour le proto-mot $ E + E * id$, une poignée peut correspondre au non terminal 
$id$ qui se réduit en $E$ en une étape grâce à la règle $ E \to id$. 

\begin{definition}[Préfixe Viable]
    Soient $G = (V, \Sigma, S, P)$ une grammaire et $w$ un proto-mot de $G$. 
    Une \emph{préfixe viable} est toute chaîne de symbole $ \alpha\in V^*$ telle que : 
        \[ \exists \beta \gamma \text{  tels que  }  S' \to^* \alpha \beta \gamma \text{ et } \beta \gamma \text{ contient une poignée complète} \] 
    où $S' \to S$ est l'axiome augmenté.   
\end{definition}

Un préfixe viable est exactement ce que l'on veut avoir au dessus de la pile d'un analyseur $LR$ si l'on ne veut 
pas de retrouver dans une impasse : 
\begin{itemize}
    \item Si la pile contient un préfixe viable, on est dans une configuration valide et l’analyse peut continuer.
    \item Si la pile contient quelque chose qui n’est pas un préfixe viable, alors l’analyseur détecte une erreur syntaxique.
\end{itemize}

\begin{definition}[Configuration Valide]
    Soient $G = (V, \Sigma, S, P)$ une grammaire. Une configuration d'un analyseur ascendant se décrit par : 
    \begin{itemize}
        \item Le \emph{contenu de la pile}
        \item Le \emph{reste de l'entrée} non consommé. 
    \end{itemize}
    Une configuration est dite \emph{valide} si la pile contient un \emph{préfixe viable} qui peut mener à 
    une dérivation complète de l'axiome. 
\end{definition}

Une configuration valide est une sorte d'instantanné de l'analyse où l'on n'a pas fait d'erreur de parsing. 

\subsection{Analyseur LR(k)}

Cette ensemble de définitions nous permet de définir formellement les grammaires $LR(k)$ de la 
façon suivante. 

\begin{definition}[Grammaire $LR(k)$]
    Soit \(G = (V, \Sigma, P, S)\) une grammaire, où :
    \begin{itemize}
        \item \(V\) est l'ensemble des symboles (terminaux et non-terminaux),
        \item \(\Sigma \subseteq V\) est l'ensemble des terminaux,
        \item \(P\) est l'ensemble des productions de la forme \(A \to \alpha\),
        \item \(S \in V \setminus \Sigma\) est l'axiome.
    \end{itemize}

    On dit que \(G\) est une \emph{grammaire LR(k)} si et seulement si,
    pour toute dérivation à droite :
    \[
    S \Rightarrow^{*} \gamma A \omega \Rightarrow \gamma \beta \omega
    \]
    avec \(A \to \beta \in P\), \(\gamma,\omega \in V^{*}\),
    il existe un \emph{unique} choix de production \(A \to \beta\) qui peut être appliqué
    pour réduire la \emph{poignée} \(\beta\),
    en ne regardant que :
    \begin{itemize}
        \item le \emph{préfixe viable} \(\gamma \beta\) (ce qui est déjà sur la pile),
        \item et au plus \(k\) symboles de lookahead dans \(\omega\).
    \end{itemize}

    En d'autres termes, la suite \(\beta\) est \textbf{une manche unique et non ambiguë}
    que l'analyseur peut identifier avec seulement \(k\) symboles d'anticipation.
\end{definition}

\begin{remark}
    On parlera identiquement de \emph{grammaire $LR(k)$} ou \emph{d'analyseur LR(k)}. 
    Le terme $LR(k)$ signifie : 
    \begin{itemize}
        \item \textbf{L}\emph{eft to right} : l'analyseur lit les mots de gauche à droite. 
        \item \textbf{R}\emph{ightmost derivation} : les motifs spécifiés par la grammaire sont recherchés par suffixes 
        de la chaîne lue. 
        \item $k$ : correspond au nombre de lettres non consommées considérées pour le choix de la dérivation à choisir. 
    \end{itemize}
\end{remark}

Définissons maintenant un automate capable de reconnaître de telles grammaires. 

\begin{definition}[Automate LR (bottom-up)]
    Un \emph{automate LR} est un \emph{automate à pile} qui reconnaît les préfixes viables d'une grammaire 
    hors contexte $ G = (V, \Sigma, S, P)$ où : 
    \begin{itemize}
        \item la \emph{pile} contient les symboles et états représentants un \emph{préfixe viable}. 
        \item l'entrée est la chaîne de symboles à analyser terminant par le caractère de fin \$. 
    \end{itemize}
    L'objectif est de lire la chaîne de gauche à droite en réduisant progressivement toutes poignées 
    dans le but d'obtenir l'axiome sur le sommet de la pile et d'avoir consommé toute l'entrée. 

    Formellement, l'automate peut être défini par : 
        $$ \boxed{ M = (Q, \Sigma, \Gamma, ACTION, GOTO, q_0, Z_0, F)} $$ 
    où : 
    \begin{itemize}
        \item $Q$ est l'ensembles des états de l'automate. 
        \item $ \Sigma$ représente l'alphabet terminal (lettres). 
        \item $ \Gamma = \Sigma \cup V_N$ est l'ensemble des symboles de la pile. 
        \item $q_0$ est l'état initial. 
        \item $Z_0$ est le symbole initial de la pile. 
        \item $F$ est l'ensemble des états acceptants. 
    \end{itemize}
    On définit en plus deux tables qui définissent la dynnamique de l'analyseur : 
    \begin{itemize}
        \item La table $ACTION$ qui détermine quelle action effectuer à la lecture d'une lettre : 
            \[ \boxed{ACTION : Q \times ( \Sigma \cup \{\$\}) \longrightarrow \{shift(q), reduce(A \to \beta), accept, error\} } \] 
        Elle gère les transitions sur les terminaux. 
        \item La table $GOTO$ qui détermine l'état dans lequel doit se trouver l'automate après 
        une réduction : 
            \[ \boxed{GOTO : Q \times V \longrightarrow Q} \] 
        Elle gère les transitions sur les non terminaux. 
    \end{itemize}
\end{definition}

\begin{definition}[Shift et Reduce]
    Comme dit précédement, nous n'avons que deux opérations à effectuer lors de la lecture d'un mot par une grammaire 
    $LR(k)$ (ou un automate LR) : \emph{shift} ou \emph{reduce}. Définissons les proprement : 
    \begin{enumerate}
        \item L'opération \emph{shift} consiste à déplacer le prochain symbole à lire sur la pile et aller 
        à l'état correspondant. On avance donc dans la lecture mais sans faire de réduction. 
        \item L'opération \emph{reduce} consiste, lorsque le sommet correspond à une poignée $ \beta$ pour 
        une règle de production $A \to \beta$, de : 
            \begin{itemize}
                \item Dépiler $ | \beta|$ symboles et états 
                \item Empiler le non-terminal $A$ 
                \item Aller dans l'état déterminé par la table $GOTO$
            \end{itemize}
        \item \emph{Accept} : si la pile est seulement constituée de $S'$ et que l'entrée est $\$$, alors le mot est reconnu. 
        \item \emph{Error} : si aucune opération shift/reduce n'est possible alors la configuration est invalide et le mot n'est pas reconnu. 
    \end{enumerate}
\end{definition}

Nous savons maintetant analyser un mot par un automate mais il reste un problème : comment définir les états 
de l'automate ? Pour cela, nous allons avoir besoin de définir le concept de \emph{fermeture} ou \emph{clôture}.

\begin{definition}[Fermeture]
    Soit $G' = \{V, \Sigma, P, S'\}$ une grammaire augmentée. 
    Un \emph{item $LR(0)$} est une règle de la forme : 
        \[ A \to \alpha \cdot \beta \] 
    où "$\cdot$" indique la position de la l'analyse dans la production $A \to \alpha \beta$. 
    La \emph{fermeture} d'un ensemble d'items $I$, notée $cl(I)$ est définie récursivement ainsi : 
    \begin{enumerate}
        \item $ \forall \text{ item} \in I, \text{ item}  \in cl(I)$. 
        \item Pour chaque item $A \to \alpha \cdot B \beta \in cl(I)$, et pour chaque production $ B \to \gamma \in P$, 
            alors $ B \to . \mu \in cl(I)$. 
    \end{enumerate}
\end{definition}

\begin{example}[Calcul de fermeture]
    Soit $G$ une grammaire dont les règles de production sont : 
        \[ 
            \begin{cases}
                S' \to S \\ 
                S \to AA \\ 
                A \to a 
            \end{cases}
        \] 
    On pose : 
        $$ I_0 := \{S' \to \cdots S\} $$ 
    Calculons maintenant les fermetures : 
    \begin{itemize}
        \item Le point est avant $S$ donc on ajoute $ S \to \cdot AA$ grâce à la règle $ S \to AA$. 
        \item Récursivement, le point est avant $A$ donc on ajoute $ A \to \cdot a$ grâce à la règle $A \to a$. 
        \item Le point $a \in \Sigma$ est un terminal on s'arrête donc ici. 
    \end{itemize}
    On a donc : 
        \[ cl(I_0) = \{S' \to \cdot S, S \to \cdot AA, A \to \cdot a\} \]
\end{example}

\begin{figure}[htbp]
    \centering
    \resizebox{\textwidth}{!}{%
    \begin{tikzpicture}[
        node distance=3cm,
        >=Stealth,
        every state/.style={minimum size=8mm, thick, draw=black!80, fill=blue!10},
        every node/.style={font=\small}
    ]
        % états (align=center permet d'utiliser \\ dans le contenu)
        \node[state, initial, align=center] (I0) {$I_0$\\$\varepsilon$};
        \node[state, right=of I0, align=center] (I1) {$I_1$\\$n$};
        \node[state, right=of I1, align=center] (I2) {$I_2$\\$n+$};
        \node[state, right=of I2, align=center] (I3) {$I_3$\\$n+n$};

        % transitions (shifts)
        \draw[->, thick] (I0) -- node[above] {$n$} (I1);
        \draw[->, thick] (I1) -- node[above] {$+$} (I2);
        \draw[->, thick] (I2) -- node[above] {$n$} (I3);

        % handles / réductions (annotation)
        \node[below=1.2cm of I1] (red1) {Handle: $n \Rightarrow E$};
        \draw[->, dashed, red] (I1) -- (red1);

        \node[below=1.2cm of I3] (red2) {Handle: $E+E \Rightarrow E$};
        \draw[->, dashed, red] (I3) -- (red2);
    \end{tikzpicture}% 
    }
    \caption{Préfixes viables et leur correspondance avec les états de l'automate LR. 
    Chaque état $I_i$ représente un ensemble de \emph{préfixes viables} atteignables. 
    Les flèches indiquent les transitions de lecture (shift). Les lignes rouges en pointillés indiquent 
    les poignées que l'on peut réduire à partir de cet état.}
\end{figure}


\begin{example}
    Travaillons autour d'un exemple pour bien saisir toutes ces nouvelles notions. 
    Soit $G$ une grammaire dont les règles de production sont : 
        \[ 
            \begin{cases}
                S' \to S \\ 
                S \to AA \\ 
                A \to a 
            \end{cases}
        \] 
    On a ajouté un symbole $S'$ fictif pour avoir un axiome de départ unique. 
    Comme vu précédement, commençons par calculer les clôtures : 
    \begin{itemize}
        \item On calcule la clôture de $I_0$ qui donne : 
            $$ cl(I_0) = \{S' \to \cdot S, S \to \cdot AA, A \to \cdot a\}$$ 
        \item Après la lecture de $S$ nous devons passer dans un nouvel état $I_1$ d'où : 
            $$ cl(I_1) = \{S' \to S \cdot\} $$ 
        \item Après la lecture de $A$ dans $I_0$ on obtient l'ensemble $I_2$ où : 
            $$ cl(I_2) = \{S \to A \cdot A, A \to \cdot a\} $$ 
        \item Après la lecture de $A$ depuis $I_2$, on obtient $I_3$ : 
            $$ cl(I_3) = \{S \to AA \cdot\} $$ 
        \item Pour finir, après la lecture de $a$ depuis $I_0$ ou $I_2$, on a : 
            $$ cl(I_4) = \{A \to \cdot a\} $$
    \end{itemize}
    On a donc les transitions suivantes : 
        $$ \text{GOTO}(I_0, S) = I_1, \quad \text{GOTO}(I_0, A) = I_2, \quad \text{GOTO}(I_0, a) = I_5 $$ 
        $$ \text{GOTO}(I_2, A) = I_4, \quad \text{GOTO}(I_2, a) = I_4 $$ 

    \begin{center}
        \begin{minipage}{0.60\textwidth}
            On peut donc construire la table action (transitions sur les terminaux) : 
            \[ 
                \begin{array}{|c|c|c|}
                    \hline 
                    \text{ État } & a & \$ \\ 
                    \hline 
                    I_0 & \text{shift } I_4 & \text{error} \\ 
                    \hline 
                    I_1 & \text{error} & \text{accept} \\ 
                    \hline 
                    I_2 & \text{shift } I_4 & \text{error} \\ 
                    \hline 
                    I_3 & \text{ reduce } S \to AA & \text{ reduce } S \to AA \\ 
                    \hline 
                    I_4 & \text{ reduce } S \to A &  \text{ reduce } S \to A \\ 
                    \hline         
                \end{array}
            \] 
        \end{minipage}
        \hfill 
        \begin{minipage}{0.25\textwidth}
             Et la table goto sur les non terminaux : 
                \[ 
                    \begin{array}{|c|c|c|}
                        \hline 
                        \text{État} & S & A \\
                        \hline 
                        I_0 & I_1 & I_2 \\ 
                        \hline 
                        I_1 & - & - \\ 
                        \hline 
                        I_2 & - & I_3 \\ 
                        \hline 
                        I_3 & - & - \\ 
                        \hline 
                        I_4 & - & - \\ 
                        \hline  
                    \end{array}
                \] 
        \end{minipage}
    \end{center}

    Résumons le traitement de l'entrée \(w = a\,a\$\) : 

    \begin{center}
        \renewcommand{\arraystretch}{1.2}
            \resizebox{\textwidth}{!}{%
            \begin{tabular}{>{\ttfamily}c >{\ttfamily}c >{\ttfamily}c l}
                \toprule
                \textbf{Pile (état/symbole)} & \textbf{Entrée} & \textbf{Action} & \textbf{Commentaire} \\
                \midrule
                $I_0$ & a a \$ & s$I_4$ & Décalage sur \texttt{a}, empiler $I_4$ \\
                $I_0$ a $I_4$ & a \$ & r($A \to a$) & Réduction : dépiler a et $I_4$, empiler A puis GOTO[$I_0$,A]=$I_2$ \\
                $I_0$ A $I_2$ & a \$ & s$I_4$ & Décalage sur \texttt{a}, empiler $I_4$ \\
                $I_0$ A $I_2$ a $I_4$ & \$ & r($A \to a$) & Réduction : dépiler a et $I_4$, empiler A puis GOTO[$I_2$,A]=$I_3$ \\
                $I_0$ A $I_2$ A $I_3$ & \$ & r($S \to A A$) & Réduction : dépiler deux $A$ et états, empiler $S$ puis GOTO[$I_0$,S]=$I_1$ \\
                $I_0$ S $I_1$ & \$ & accept & Mot reconnu, analyse terminée \\
                \bottomrule
            \end{tabular}%
            }
    \end{center}
\end{example}


